%% sec1_intro.tex -- Section 1: Introduction

\section{Introduction}
\label{sec:intro}

A $(v,k)$-lottery is a game in which $k$ distinct numbers are drawn
uniformly at random from $[v] = \{1,\ldots,v\}$.  A player holding a ticket,
which is a $k$-subset of $[v]$, wins a prize when at least $t$ of the drawn
numbers appear on the ticket.  Games of this form are played in many countries
and have attracted sustained mathematical interest since the work of
Sch\"{o}nheim \cite{Schoenheim1964}, who introduced the study of covering
numbers $C(v,k,t)$, the minimum cardinality of a family of $k$-subsets
of $[v]$ such that every $t$-subset is contained in at least one.

The bulk of the literature pursues a deterministic goal: given $n$ tickets,
arrange them so that every possible draw $R$ satisfies $|A_i \cap R| \ge t$
for at least one $A_i$.  This is the $(v,k,t,m)$-lotto design framework
formalised in \cite{Furedi1996} and tabulated by Li and van Rees \cite{Li2002},
who describe greedy and simulated annealing constructions for small instances.
Gordon, Kuperberg and Patashnik \cite{Gordon1995} give efficient upper-bound
constructions; comprehensive tables appear in the \emph{Handbook of
Combinatorial Designs} \cite{Colbourn2007}.

The deterministic paradigm has a clear virtue: a well-chosen set of tickets
guarantees a prize regardless of the draw.  For practical budgets, however,
the number of tickets needed grows quickly with $v$ and $t$, and a player
committing to $n$ tickets faces a different question: given exactly this many
tickets, which selection maximises the probability of winning at least once?

This probabilistic question is the subject of the present paper.  Let $R$ be
a uniformly random $k$-subset of $[v]$ and let $\mathcal{A} = (A_1,\ldots,A_n)$
be a collection of $n$ tickets.  We study
\begin{equation}
  \label{eq:objective}
  P_{\mathrm{win}}(\mathcal{A})
  = P\!\left(\,\bigcup_{i=1}^{n} \{|A_i \cap R| \ge t\}\right).
\end{equation}
Because the marginal prize probability $p(v,k,t) = P(|A_i \cap R| \ge t)$ is
the same for every ticket (it depends only on $v$, $k$, $t$), the player's
only lever is the joint distribution of the scores, which is controlled by
the pairwise overlaps $|A_i \cap A_j|$.

The central finding of the paper is that minimising pairwise overlap
maximises $P_{\mathrm{win}}$.  For two tickets we prove this rigorously
(Theorem~\ref{thm:two-ticket-monotone}).  A complementary calculation shows
that for the Lotof\'{a}cil parameters $(v=25,k=15,t=11)$, any two tickets
with overlap at most six have mutually exclusive winning events, so
$P_{\mathrm{win}} = 2p$ exactly (Proposition~\ref{prop:mutual-excl}).  This
is the maximum possible for $n=2$ and follows purely from the arithmetic of
the draw.  For $n \ge 3$ we state a monotonicity conjecture backed by
simulation, and show that two complementary algorithms produce near-optimal
collections efficiently.

The first algorithm, \emph{Greedy Coverage Selection}, adds tickets one at a
time by maximising the number of new universe elements covered, with ties
broken by cumulative overlap.  It inherits the $(1-1/e)$-approximation
guarantee of \citet{Nemhauser1978} for the maximum-coverage objective.  The
second algorithm applies \emph{Simulated Annealing} to the greedy solution,
directly minimising the total overlap $\sum_{i<j}|A_i \cap A_j|$ via
single-element swaps \cite{Kirkpatrick1983}.

To validate both algorithms we study the Brazilian Lotof\'{a}cil lottery,
which, to our knowledge, appears in the academic literature for the first
time here.  The lottery's high ticket-size-to-universe ratio ($k/v = 15/25$)
forces non-trivial pairwise overlap and places it in a qualitatively
different regime from the 6/49-type games that have been studied theoretically.
Using a paired Monte Carlo design with $50{,}000$ shared simulated draws per
condition, we estimate $P_{\mathrm{win}}$ for $n$ ranging from 1 to 100.
The results show gains of up to 11 percentage points over random selection
at $n=5$ (paired $t$-test: $t = 18.66$, $p < 10^{-23}$, Cohen's $d = 2.64$),
with statistically significant advantages at every $n \ge 2$ up to the
saturation point near $n = 50$.

The closest precursor in the algorithmic literature is \citet{Li2002}, who
tabulate optimal lotto designs using greedy algorithms and simulated annealing,
but with the deterministic coverage objective.  \citet{Mohammadi2012} formulate
lottery selection as an integer program and propose a meta-heuristic for small
instances.  A recent working paper by \citet{Liu2024} studies a probabilistic
formulation for the UK 6/59 lottery and finds that balanced pairwise overlap
maximises expected return across all prize tiers; our Theorem~\ref{thm:two-ticket-monotone},
which concerns the top prize level only, is consistent with that finding.
None of these works address the high-$k$ regime of the Lotof\'{a}cil, provide
a full statistical characterisation across a range of $n$, or prove the
mutual-exclusivity result of Proposition~\ref{prop:mutual-excl}.

The paper is organised as follows.  Section~\ref{sec:framework} establishes
notation and the basic probabilistic model.
Section~\ref{sec:theory} develops the overlap monotonicity results and
degenerate cases.  Section~\ref{sec:algorithms} describes the algorithms and
their guarantees.  Section~\ref{sec:empirical} presents the Lotof\'{a}cil
experiment.  Section~\ref{sec:stats} provides the statistical analysis.
Section~\ref{sec:discussion} discusses connections to covering theory,
extensions, and limitations.
